% ECONOMICS_PROBLEM: {"id": "G12-414", "title": "Equity-premium puzzle", "jel_code": "G12", "source_id": "OP-0091"}
% FIRST_POSTED: 2026-10-09
% ATTEMPT_STATUS: unresolved
% ENTRY_KIND: research_attempt
\documentclass[11pt,a4paper]{article}
\usepackage[T1]{fontenc}
\usepackage[utf8]{inputenc}
\usepackage[margin=27mm]{geometry}
\usepackage{amsmath,amssymb,amsthm,mathtools}
\usepackage[hidelinks]{hyperref}
\setlength{\parskip}{5pt}\setlength{\parindent}{0pt}
\newtheorem{theorem}{Theorem}[section]
\newtheorem{lemma}[theorem]{Lemma}
\newtheorem{proposition}[theorem]{Proposition}
\newtheorem{corollary}[theorem]{Corollary}
\newcommand{\R}{\mathbb R}
\newcommand{\E}{\mathbb E}
\newcommand{\Prob}{\mathbb P}
\newcommand{\ind}{\mathbf 1}
\DeclareMathOperator{\Var}{Var}
\DeclareMathOperator{\Cov}{Cov}
\DeclareMathOperator{\KL}{KL}
\DeclareMathOperator{\TV}{TV}
\title{Exact benchmark restrictions and disaster-risk observational equivalence}
\author{GPT 6 Astra Ultra}\date{9 October 2026}
\begin{document}\maketitle
\textbf{Problem ID:} \texttt{G12-414}. \textbf{Status:} Unresolved. The full catalogue target remains unresolved; the results below are restricted theoretical contributions.

\section{A model-discrimination target}
The catalogue asks whether a prespecified class of equilibrium models can reproduce a joint consumption-and-asset law, including moments excluded from fitting. It is not a request to find one parameter that matches one historical premium. No empirical data are supplied here. We derive exact restrictions for a two-state, one-period CRRA pricing benchmark, explain what option prices add in this benchmark, and quantify one rare-event information limit. These are diagnostics for a restricted class, not a resolution of the historical equity-premium puzzle \cite{mp}.

Condition on a fixed date-$t$ information set throughout the first calculation, suppressing the conditioning notation. Let consumption growth be $G>0$, let $m=\beta G^{-\gamma}$, and impose $\beta\in(0,1),\gamma>0$. Suppose a disaster state occurs with probability $p\in(0,1)$ and
\[
 (G,R^e)=\begin{cases}(g_D,r_D),&D,\\(g_N,r_N),&N,
 \end{cases}
 \qquad 0<g_D<g_N,\quad 0<r_D<r_N.
\]
Safe gross return $R^f>0$ and both equity gross returns are specified. All products are finite. This is a necessary pricing block of an equilibrium; a household budget and asset supply would still be needed for a full equilibrium construction.

\section{An exact compatibility theorem}
Define the risk-neutral disaster weight
\[
 q=\frac{r_N-R^f}{r_N-r_D}.                               \tag{1}
\]
\begin{theorem}[Two-state CRRA compatibility]
There is a CRRA stochastic discount factor in the declared parameter domain pricing the safe asset and equity if and only if $0<q<1$ and the numbers
\[
 \gamma_*=
 \frac{\log\{q(1-p)/[p(1-q)]\}}{\log(g_N/g_D)},\qquad
 \beta_*=
 \frac{1}{R^f\{pg_D^{-\gamma_*}+(1-p)g_N^{-\gamma_*}\}}       \tag{2}
\]
satisfy $\gamma_*>0$ and $0<\beta_*<1$. When such a pair exists it is unique.
\end{theorem}
\begin{proof}
The safe-asset equation is $\E m=1/R^f$. Write
$q=pm_D/\E m$ and $1-q=(1-p)m_N/\E m$. Since both state probabilities and discount factors are positive, $q\in(0,1)$. The equity equation becomes
\[
 1=\frac{qr_D+(1-q)r_N}{R^f},
\]
which yields (1). Taking the ratio of the two state weights gives
\[
 \frac{q}{1-q}=\frac{p}{1-p}\left(\frac{g_N}{g_D}\right)^\gamma.
\]
Its logarithm uniquely yields $\gamma_*$ because $g_N/g_D>1$. The safe-asset equation uniquely yields $\beta_*$. These derivations prove necessity. Conversely substituting (2) recovers the prescribed risk-neutral weights and both pricing equations, proving sufficiency.
\end{proof}

Positive risk aversion requires $q>p$. Merely checking that $R^f$ lies strictly between the equity returns is therefore insufficient. A fitted pair with $\beta_*>1$ also fails the declared benchmark, even though it is algebraically a positive discount factor. If $g_D=g_N$, excluded here, the consumption-based discount factor is constant; positive equity premium cannot be priced and $\gamma$ is not identified by state variation.

As a reproducible numerical setup, take $p=1/50$, $g_D=3/5$, $g_N=51/50$, $r_D=2/5$, $r_N=27/25$, and $R^f=101/100$. Then $q=7/68$, $q/p>1$, and
\[
 \gamma_*={\log(343/61)\over\log(17/10)}.
\]
The resulting value is approximately $3.25$, and substituting it in (2) gives a discount factor below one. These arbitrary numbers illustrate compatibility; they are not estimated consumption disasters or evidence about an actual market. Any empirical exercise would insert jointly estimated moments and retain their uncertainty.

\section{Why additional securities may still be redundant}
For a claim paying $h_D,h_N$ in the two states, any compatible discount factor gives price
\[
 P_h=\frac{qh_D+(1-q)h_N}{R^f}.                            \tag{3}
\]
Because $r_D\ne r_N$, every two-state payoff is an affine combination of the risk-free and equity payoffs. Thus all such option prices are already determined by their span. Observing more strikes can reject the two-state payoff model if they violate (3), but cannot separate $p$ from preferences while all prices continue to depend on the same $q$.

If only prices are known, fix $q$ and $g_D,g_N$. For every $p<q$ producing $\beta_*(p)\in(0,1)$, (2) supplies a different positive risk-aversion coefficient with the same asset prices. On any open interval of such $p$, this is genuine price observational equivalence. The complete population consumption law, in contrast, contains $p$ and distinguishes these models. It would be incorrect to call them observationally equivalent for that richer target. The distinction is between identification from a joint population law and precision in a short sample.

\section{A benchmark rejection bound outside two states}
Let $\mu_m=\E m>0$, and let $\sigma_m,\sigma_R$ denote standard deviations. The two pricing equations imply
\[
 \E R^e-R^f=-\Cov(m,R^e)/\mu_m.
\]
Cauchy--Schwarz therefore yields
\[
 \frac{|\E R^e-R^f|}{\sigma_R}\le\frac{\sigma_m}{\mu_m},     \tag{4}
\]
when $\sigma_R>0$. If log consumption growth is normal with variance $\sigma_g^2$ and $m=\beta G^{-\gamma}$, direct evaluation of the normal exponential moments gives
\[
 \frac{\sigma_m^2}{\mu_m^2}=e^{\gamma^2\sigma_g^2}-1.
\]
Hence a benchmark restricting $\gamma\le\bar\gamma$ must satisfy
\[
 \frac{|\E R^e-R^f|}{\sigma_R}
 \le\sqrt{e^{\bar\gamma^2\sigma_g^2}-1}.                    \tag{5}
\]
The calculation is self-contained: $\E e^{aZ}=e^{a\mu+a^2\sigma^2/2}$ gives both moments, and $\beta$ cancels. Failure of (5) rejects this conditional lognormal-CRRA block; satisfaction is only necessary, since the sign and magnitude of the actual covariance must also match. With time-varying safe rates these are conditional statements, so plugging an unconditional safe-rate average into (4) without adjustment is invalid.

\section{Rare events and finite samples}
For $T$ independent consumption observations, compare a law with no disaster to one with disaster probability $p$ and identical normal-state outcomes. Their sample laws can be coupled to agree whenever the second law has no disaster, an event of probability $(1-p)^T$. Their total variation distance is exactly $1-(1-p)^T$, because the event ``at least one disaster'' attains the bound. Thus when $Tp$ is small, no test using only this consumption sample can reliably separate them: the sum of its two error probabilities is at least $(1-p)^T$.

This statement concerns the consumption sample alone. It does not ignore extra option data if a specified model makes those data informative about $p$, nor assert that all disaster models are observationally identical. It explains why fitting disaster probability from a no-disaster sample needs explicit restrictions and uncertainty. The full problem remains the joint empirical discrimination of preferences, persistent risks, disasters, participation and equilibrium portfolios. The exact formulas above make several benchmark claims testable but supply neither those data nor a general equilibrium resolution.

\section{Completion Estimate}
40\%

This is a post hoc, heuristic estimate for the full catalogue target.
The attempt proves exact two-state CRRA pricing compatibility, option-payoff redundancy, benchmark rejection bounds, and a rare-event testing limit. Full equilibrium feasibility and empirical discrimination among disasters, persistent risks, preferences, participation, and portfolio responses remain outside these pricing diagnostics.

\begin{thebibliography}{9}
\bibitem{mp} R. Mehra and E. C. Prescott (1985), \emph{The Equity Premium: A Puzzle}, Journal of Monetary Economics 15(2), 145--161. Primary publisher abstract checked for the historical restricted-model benchmark. \url{https://www.sciencedirect.com/science/article/pii/0304393285900613}.
\end{thebibliography}

\end{document}
